\chapter{Stochastic Gradient Descent}
Assume we want to solve
\begin{equation}
    \min_x f(x),
\end{equation}
but we do have access only to random but unbiased estimates of $f$, respectively $\partial f$. That is, we assume we can evaluate $G(x,z)$ which satisfies
\begin{equation}\label{stochastic:eq:gradient}
    \E[G(x,Z)] \in \partial f(x)
\end{equation}
where $Z$ is some random variable.
It turns out that many of the convergence proofs can be transferred to this stochastic setting with almost no extra effort. Before showing this, let us motivate this setting.

In modern machine learning, most learning problems can be formulated as
\begin{equation}
    \min_x f(x)\coloneqq \sum_{i=1}^N f_i(x)
\end{equation}
very $N$ is usually the number of training samples. For instance, if we want to learn a parametrized map $f_\theta$ which should approximate some relation $x\mapsto y$ based on the training data $(x_i,y_i)_{i=1}^N$, the corresponding training problem could be
\begin{equation}
    \min_\theta \sum_{i=1}^N \|y_i-f_\theta(x_i)\|^2.
\end{equation}
Due to the encountered sizes of training data, it is often not possible to directly compute gradients of $f$ because of memory limitations.
If $f_i$ are differentiable a possible remedy is the stochastic update
\begin{equation}
    \begin{cases}
        \text{Choose $i\in\Ic$}\\
        x_{k+1} = x_k - \tau \nabla f_i(x_k)
    \end{cases}
\end{equation}
where $\Ic=\{1,\dots,N\}$.
Defining $Z_k$ as a sequence of iid uniformly distributed random variables in $\Ic$ the above can be written as
\begin{equation}
    x_{k+1} = x_k - \tau G(x_k,Z_k)
\end{equation}
with $G(x,z) \coloneqq \nabla f_z(x)$. One can easily check that such $G$ satisfies~\eqref{stochastic:eq:gradient}.

\begin{theorem}
    Let $f$ be $L$-Lipschitz and assume the step sizes satisfy
    \[
        \frac{\sum_{k=1}^n\tau_k}{\sum_{k=1}^n\tau_k^2}\rightarrow \infty
    \]
    as $n\rightarrow\infty$. 
    Moreover, assume the stochastic gradients are unbiased and of bounded variance, specifically, for all $x\in \R^d$ we have
    \begin{equation}
        \begin{cases}
            \E[G(x,Z)] \in &\partial f(x)\\
            \E[\|G(x,Z)\|^2] - \|\E[G(x,Z)]\|^2 &\leq \sigma^2<\infty
        \end{cases}
    \end{equation}
Then it holds $\E[f^n_{\mathrm{best}}-f(x^*)]\rightarrow 0$.
\end{theorem}
\begin{proof}
    \begin{equation}
        \begin{aligned}
            \|x_{k+1}-x^*\|^2
            &= \|\proj_C(x_{k}-\tau_k G(x_k,Z_k))-\proj_C(x^*)\|^2\\
            &\leq \|x_{k}-\tau_k G(x_k,Z_k)-x^*\|^2\\
            &= \|x_{k}-x^*\|^2 - 2\tau_k\inner{x_k-x^*}{G(x_k,Z_k)} + \tau_k^2\|G(x_k,Z_k)\|^2.
        \end{aligned}
    \end{equation}
    Note that
    \begin{equation}
        \begin{aligned}
            \E[\inner{x_k-x^*}{G(x_k,Z_k)}]
            &= \E[\E[\inner{x_k-x^*}{G(x_k,Z_k)}|x_k]]\\
            &= \E[\inner{x_k-x^*}{\E[G(x_k,Z_k)|x_k]}]\\
            &\geq \E[f(x_k) - f(x^*)]\\
            &= \E[f(x_k)] - f(x^*)
        \end{aligned}
    \end{equation}
    as well as
    \begin{equation}
        \E[\|G(x_k,Z_k)\|^2] 
        = \E[\|G(x_k,Z_k)\|^2] - \|E[G(x_k,Z_k)]\|^2 + \|E[G(x_k,Z_k)]\|^2
        \leq \sigma^2 + L^2
    \end{equation}
    Thus, taking expectation leads to 
    \begin{equation}
        \begin{aligned}
            \E[\|x_{k+1}-x^*\|^2]
            \leq \E[\|x_{k}-x^*\|^2] - 2\tau_k (\E[f(x_k)] - f(x^*)) + \tau_k^2(\sigma^2 + L^2).
        \end{aligned}
    \end{equation}
    The remaining proof is as in the deterministic case. Summing over $k$ leads to
    \begin{equation}
        \begin{aligned}
            \sum_{k=0}^{K-1}\tau_k (\E[f(x_k)] - f(x^*))\leq \frac{1}{2}\E[\|x_{0}-x^*\|^2] + (\sigma^2 + L^2)\sum_{k=0}^{K-1}\tau_k^2
        \end{aligned}
    \end{equation}
    Defining
    \begin{equation}
        \sigma_n = \frac{\sum_{k=0}^{K-1}\tau_k}{\sum_{k=0}^{K-1}\tau_k^2}
    \end{equation}
    we obtain
    \begin{equation}
        \begin{aligned}
            \sigma_n \E[f^n_{\mathrm{best}}-f(x^*)]
            =& \frac{\sum_{k=0}^{K-1}\tau_k\E[f^n_{\mathrm{best}}-f(x^*)]}{\sum_{k=0}^{K-1}\tau_k^2}\\
            \leq& \frac{\sum_{k=0}^{K-1}\tau_k\E[f(x_k)-f(x^*)]}{\sum_{k=0}^{K-1}\tau_k^2}\\
            \leq& \frac{\E[\|x_{0}-x^*\|^2]}{2\sum_{k=0}^{K-1}\tau_k^2} + \frac{(\sigma^2 + L^2)\sum_{k=0}^{K-1}\tau_k^2}{\sum_{k=0}^{K-1}\tau_k^2}\\
        \end{aligned}
    \end{equation}
    and thus
    \begin{equation}
        \E[f^n_{\mathrm{best}}-f(x^*)] \leq \frac{\E[\|x_{0}-x^*\|^2]}{2\sum_{k=0}^{K-1}\tau_k} + \frac{(\sigma^2 + L^2)\sum_{k=0}^{K-1}\tau_k^2}{\sigma_n}
    \end{equation}
    which goes to zero as $n\rightarrow\infty$.
\end{proof}